% Section content template for individual Educative lessons
% This file contains the actual content for a specific section
% Generated from Educative API section content

% Section: Code for the Amazon Online Shopping System
% Section ID: 5634660812193792
% Section Slug: code-for-the-amazon-online-shopping-system
% Generated: 2025-12-12T22:46:06.576498

We've reviewed Amazon's different aspects and observed the attributes attached to the problem using various UML diagrams. Now , let's explore the more practical side of things, where we'll implement Amazon using multiple languages. This is usually the last step in an object-oriented design interview process.
\\
We have chosen the following languages to write the skeleton code of the different classes present in Amazon:

\begin{itemize}
\item
\protect\phantomsection\label{sGareDrLIyf6WSi7gZREL}
Java
\item
\protect\phantomsection\label{veQRBPlFazZNvZpQDR6iZ}
C\#
\item
\protect\phantomsection\label{VNLCRVe6vx12AYoq4GhBX}
Python
\item
\protect\phantomsection\label{CIPLXwc9Gld31YlE7NrGB}
C++
\item
\protect\phantomsection\label{ZF8PWbh_JW__I7JBSsSGy}
JavaScript
\end{itemize}
\\
If you want to skip directly to the implementation and see the complete workflow, simply scroll down or click~\hyperref[Executable-code-Amazon-online-shopping-system]{Executable Code: Amazon Online Shopping System}.

\subsection{Amazon classes}\label{3Kc_WSGNh9BOYfaEPXewu}
\\
This section will provide the skeleton code of the classes designed in the class diagram lesson.
\begin{quote}
\textbf{Note:} For simplicity, we are not defining getter and setter functions. The reader can assume that all class attributes are private, accessed through their respective public getter methods, and modified only through their public method functions.
\end{quote}
\subsubsection{Constants}\label{1DftBh-vIsuKI400oYwWO}
\\
The following code defines the various enums and custom data types being used in the Amazon design:
\begin{quote}
\textbf{Note:} JavaScript does not support enumerations, so we will use the Object.freeze()\ method as an alternative that freezes an object and prevents further modifications.
\end{quote}

\textbf{Constant definitions}

\begin{lstlisting}[language=Java]
public class Address {
  private int zipCode;
  private String address;
  private String city;
  private String state;
  private String country;
  
  public String getFullAddress();
}

enum OrderStatus {
  UNSHIPPED, 
  PENDING, 
  SHIPPED, 
  CONFIRMED, 
  CANCELED, 
  REFUNDED
}

enum AccountStatus {
  ACTIVE, 
  INACTIVE, 
  BLOCKED
}

enum ShipmentStatus {
  PENDING, 
  SHIPPED, 
  DELIVERED, 
  ON_HOLD
}

enum PaymentStatus {
  CONFIRMED, 
  DECLINED, 
  PENDING, 
  REFUNDED
}
\end{lstlisting}

\subsubsection{Account}\label{CaSIBJK3fyQtURMYYmpVL}
\\
The \texttt{Account} class refers to an account of a customer on Amazon and contains a customer's personal details, such as their name, shipping address, credit card information, etc. It also provides authenticated users and admins with the functionality to add products, product reviews, and reset passwords. The definition of this class is given below:

\textbf{The Account class}

\begin{lstlisting}[language=Java]
public class Account {
  private String userName;
  private String password;
  private String name;
  private List<Address> shippingAddress;
  private AccountStatus status;
  private String email;
  private String phone;

  private List<CreditCard> creditCards;
  private List<ElectronicBankTransfer> bankAccounts;

  public Address getShippingAddress();
  public boolean addProduct(Product product);
  public boolean addProductReview(ProductReview review, Product product);
  public boolean deleteProduct(Product product);
  public boolean deleteProductReview(ProductReview review, Product product);
  public boolean resetPassword();
}
\end{lstlisting}

\subsubsection{Admin}\label{a-GjKp6-VT6SMnzmDC3jh}
\\
The \texttt{Admin} class refers to a person from the administration of Amazon who can block users, add, modify, or delete product categories. The definition of this class is provided below:

\textbf{Admin class}

\begin{lstlisting}[language=Java]
public class Admin {
  private Account account;
  public boolean blockUser(Account account);
  public boolean addNewProductCategory(ProductCategory category);
  public boolean modifyProductCategory(ProductCategory category);
  public boolean deleteProductCategory(ProductCategory category);
}
\end{lstlisting}

\subsubsection{Customer}\label{KcbwE9F21HX5kQsBlPvQF}
\\
The \texttt{Customer} class is an abstract class with the \texttt{AuthenticatedUser} and \texttt{Guest} classes derived from it. The classes are defined below:

\textbf{Customer and its child classes}

\begin{lstlisting}[language=Java]
// Customer is an abstract class
public abstract class Customer {
  private ShoppingCart cart;
  
  public abstract List<Product> searchproduct(String name);
  
  
}

public class Guest extends Customer {
  public boolean registerAccount();
  public abstract List<Product> searchproduct(String name);
}

public class AuthenticatedUser extends Customer {
  private Account account;
  private Order order;
  
  public OrderStatus placeOrder(Order order);
  public abstract List<Product> searchproduct(String name);
  public ShoppingCart getShoppingCart();
}
\end{lstlisting}

\subsubsection{Product, product category, and product review}\label{BVr3iVSKEF-uu4yWhpAAa}
\\
The following code defines the classes relating to a particular product:

\textbf{The Product, ProductCategory, and ProductReview classes}

\begin{lstlisting}[language=Java]
public class Product {
  private String productId;
  private String name;
  private String description;
  private byte[] image;
  private double price;
  private ProductCategory category;
  private List<ProductReview> reviews;
  private int availableItemCount;
  private Account account;

  public int getAvailableCount();
  public int updateAvailableCount();
  public boolean updatePrice(double newPrice);
}

public class ProductCategory {
  private String name;
  private String description;
  private List<Product> products;
}

public class ProductReview {
  private int rating;
  private String review;
  private byte[] image;
  private AuthenticatedUser user;
}
\end{lstlisting}

\subsubsection{Shopping cart and cart items}\label{3VyJ6YosyGRVL7Tb2VQ2W}
\\
The \texttt{ShoppingCart} and CartItem\ classes add items to the shopping cart, update the quantities, and send the list of items to checkout. Both classes are defined below:

\textbf{The ShoppingCart and CartItem classes}

\begin{lstlisting}[language=Java]
public class CartItem {
  private int quantity;
  private double price;
  public boolean updateQuantity(int quantity);
}

public class ShoppingCart {
  private int totalPrice;
  private List<Items> items;

  public boolean addItem(Item item);
  public boolean removeItem(Item item);
  public List<Item> getItems();
  public boolean checkout();
  public boolean verify();
  public double getTotalPrice();
}
\end{lstlisting}

\subsubsection{Order}\label{QVGf7_C5skQ4iD5RGmVI4}
\\
The \texttt{Order} class is responsible for making the payment, updating the order status, and sending the particular order for shipment. The definition of the class is given below:

\textbf{The Order class}

\begin{lstlisting}[language=Java]
import java.util.ArrayList;
import java.util.Date;
import java.util.List;

public class Order {
  private String orderNumber;
  private AuthenticatedUser cutomer;
  private Date orderDate;
  private OrderStatus status;
  private ShoppingCart shoppingCart;

  public Order(String orderNumber, OrderStatus status) {
    this.orderNumber = orderNumber;
    this.orderDate = new Date();
    this.status = status;
    
  }

  public boolean sendForShipment();
  public PaymentStatus makePayment(Payment payment);
  public boolean verify(CartItem item);
}
\end{lstlisting}

\subsubsection{Shipment}\label{t2bYEUsocE9svZycmtH4B}
\\
The \texttt{Shipment} class keeps track of all the major details of the order's dispatch and creates the shipment record. The class is defined below:

\textbf{The Shipment class}

\begin{lstlisting}[language=Java]
public class Shipment {
  private String shipmentNumber;
  private Date shipmentDate;
  private Date estimatedArrival;
  private String shipmentMethod;
  private ShipmentStatus status;

  public Shipment(String shipmentNumber, String shipmentMethod, Date estimatedArrival, ShipmentStatus status) {}
}
\end{lstlisting}

\subsubsection{Payment}\label{INuRmh5_DxMq2iIJNxzEq}
\\
The \texttt{Payment} class is another abstract class with the \texttt{ElectronicBankTranfer}, \texttt{CreditCard}, and \texttt{Cash} classes as its child classes. This takes the \texttt{PaymentStatus} enum to keep track of the payment status. The definition of this class is provided below:

\textbf{Payment and its child classes}

\begin{lstlisting}[language=Java]
// Payment is an abstract class
public abstract class Payment {
  private double amount;
  private Date timestamp;
  private PaymentStatus status;
  public abstract PaymentStatus makePayment();
}

public class CreditCard extends Payment {
  private String nameOnCard;
  private String cardNumber;
  private String billingAddress;
  private int code;
  public PaymentStatus makePayment();
}

public class ElectronicBankTransfer extends Payment {
  private String bankName;
  private String routingNumber; 
  private String accountNumber;
  private String billingAddress;
  public PaymentStatus makePayment();
}

public class Cash extends Payment {
  private String billingAddress;
  public PaymentStatus makePayment();
}
\end{lstlisting}

\subsubsection{Notification}\label{0ZVCay7W6Bkab1KXP_cLz}
\\
The \texttt{Notification} class is responsible for sending notifications to customers about the order and shipment status. It is defined below:

\textbf{Notification and its child classes}

\begin{lstlisting}[language=Java]
// Notification class
public class Notification {
  private int notificationId;
  private Date createdOn;
  private String content;

  public abstract boolean sendNotification(Account account);
}
\end{lstlisting}

\subsubsection{Search}\label{julu7b82zRwsJjYIajaE4}
\\
The \texttt{Search} class enables the search functionality based on the given criteria (name and category of product).

\textbf{The Search class}

\begin{lstlisting}[language=Java]
public class  Search {
  private HashMap<String, List<Product>> products = new HashMap<>();

  public void setProducts(Map<String, List<Product>> products);
  public List<Product> searchProductsByName(String name);
  public List<Product> searchProductsByCategory(String category);
  public Product getProductDetails(String productId);
}
\end{lstlisting}

\subsection{Executable code: Amazon online shopping system}\label{mcg2gLg9H3CiR7Ox5y7Ig}
\\
Below is a fully self-contained, runnable program in Java, C\#, C++, Python, and JavaScript that demonstrates the core workflows of the Amazon Online Shopping System. The system's main driver code resides in \texttt{Driver.java},~\texttt{Driver.cs},~\texttt{Driver.py},~\texttt{Driver.cpp}, and~\texttt{Driver.js}~for all respective languages. You can click the ``Run'' button to execute the codes.

\subsubsection{What does this code show}\label{P3iJfzgPuLN02tbkrblED}

\begin{itemize}
\item
\protect\phantomsection\label{FyQrEJX4QaUFdPLcIrze6}
\textbf{System initialization:} Initializes a guest user who browses the platform and registers an account. It also initializes an authenticated user (John) and creates a shopping cart and products to simulate the shopping experience.
\item
\protect\phantomsection\label{qPYj-qtTHhvu9YAUSn0vv}
\textbf{Scenario 1 (Guest browses and registers):} A guest user browses the platform and registers a new account with a username, password, and email. This simulates a typical scenario where new users create an account to start shopping.
\item
\protect\phantomsection\label{YmFHnOP6p2z1Uvw8965uu}
\textbf{Scenario 2 (Authenticated user adds items to cart and places an order):} The authenticated user, John, browses the available products, adds items (a laptop and a mouse) to his shopping cart, and proceeds to order. The cart is verified, and the order is placed for the first product in the cart, simulating the purchase process.
\item
\protect\phantomsection\label{j0VFEwPBoIv-Y4ZpSh8XG}
\textbf{Scenario 3 (Order payment with credit card):} John makes the payment using a credit card for the order in Scenario 2. The payment is processed, and the shipment is prepared. The shipment details and payment status are displayed, showing the final transaction completion stages.
\end{itemize}
\begin{quote}
\textbf{Hint: Try customizing the code!}\\
This code is designed for experimentation, not just reading. You can edit, extend, or modify any part of the example.
\end{quote}

\textbf{Amazon Online Shopping System}

\begin{lstlisting}[language=Java]
import java.util.*;

public class Driver {
  public static void main(String[] args) {
    System.out.println("========== Welcome to the Amazon Online Shopping System ==========
");

    // Scenario 1: Guest Browses and Registers
    System.out.println("==== Scenario 1: Guest Browses and Registers ====");
    Guest guest = new Guest();
    System.out.println("🧑 Guest browses the platform...");
    guest.registerAccount("john_doe", "password123", "john@example.com");
    
    // Scenario 2: Authenticated User Adds Items to Cart and Places an Order
    System.out.println("
==== Scenario 2: Authenticated User Adds Items to Cart and Places an Order ====");
    
    // 2.1. Authenticated User Setup
    Account johnAccount = new Account();
    johnAccount.setStatus(AccountStatus.ACTIVE);
    AuthenticatedUser john = new AuthenticatedUser();
    
    // 2.2. Create Product Catalog
    Search search = new Search();
    ProductCategory electronics = new ProductCategory("Electronics", "Devices and gadgets");
    Product laptop = new Product("P001", "Gaming Laptop", "Powerful 16GB RAM laptop", 1500.00, electronics);
    Product mouse = new Product("P002", "Wireless Mouse", "Ergonomic wireless mouse", 25.00, electronics);
    
    // Manually populate product catalog
    List<Product> electronicsList = new ArrayList<>();
    electronicsList.add(laptop);
    electronicsList.add(mouse);
    Map<String, List<Product>> productMap = new HashMap<>();
    productMap.put("electronics", electronicsList);
    search.setProducts(productMap);

    System.out.println("📦 Products available in 'Electronics':");
    for (Product p : search.searchProductsByCategory("electronics")) {
      System.out.println(" - " + p.getName() + ": $" + p.getPrice());
    }

    // 2.3. User adds items to shopping cart
    System.out.println("
🛒 User adds items to shopping cart:");
    john.addItem(laptop, 1);
    john.addItem(mouse, 2);

    // 2.4. Show cart summary
    System.out.println("
🧾 Shopping Cart Summary:");
    ShoppingCart cart = john.getShoppingCart();
    for (CartItem item : cart.getItems()) {
      System.out.println(" - " + item.getProduct().getName() + " x " + item.getQuantity() + " = $" + item.getPrice());
    }
    System.out.println("Total: $" + cart.getTotalPrice());

    // 2.5. Verify cart and place order
    if (cart.verify()) {
      System.out.println("
✅ Cart verified. Proceeding to place order...");
    }
    CartItem itemToOrder = cart.getItems().get(0);
    OrderStatus status = john.placeOrder(itemToOrder, itemToOrder.getPrice());

    // Scenario 3: Order Payment with Credit Card
    System.out.println("
==== Scenario 3: Order Payment with Credit Card ====");
    // 3.1. Process payment using CreditCard
    Payment payment = new CreditCard(itemToOrder.getPrice(), "John Doe", "1234567812345678", "123 Street, NY", 123);

    // 3.2. Order and Shipment
    Order order = new Order("ORD123", status);
    PaymentStatus paymentStatus = order.makePayment(payment);

    Shipment shipment = new Shipment("SHIP123", "FedEx", new Date(System.currentTimeMillis() + 3 * 86400000), ShipmentStatus.PENDING);

    // 3.3. Display Shipment and Payment Status
    // System.out.println("
📦 Shipment Details:");
    // for (ShipmentLog log : shipment.getShipmentLogs()) {
    //   System.out.println(" - Status: " + log.getStatus() + ", Logged on: " + log.getCreationDate());
    // }

    System.out.println("
💳 Payment Status: " + paymentStatus);
    System.out.println("📑 Order Status: " + status);

    System.out.println("
🎉 Thank you for shopping with us!");
    System.out.println("=====================================================");
  }
}
\end{lstlisting}

\subsection{Wrapping up}\label{d6DcwCJ0jKUAiX5iUpOvw-2}
\\
In this chapter, we've explored Amazon's complete design, including how it is visualized using various UML diagrams and designed using object-oriented principles and design patterns.

% End of section content